
\chapter{Glossary}
\begin{description}\small
  \item[ABI] application binary interface: how compiled code calls other compiled code. CPython's
        changes between minor versions.
  \item[auditwheel, delocate, delvewheel] tools that copy shared libraries into a wheel and fix its
        tag (Linux, macOS, Windows).
  \item[install component] a label on CMake install rules, to install a subset of them.
  \item[build backend / frontend] the backend (scikit-build-core) turns a project into a wheel; the
        frontend (pip, \code{build}, \code{uv build}) calls it in an isolated environment.
  \item[deployment target] the oldest macOS a binary supports.
  \item[editable install] an install that points to the source folder instead of copying it.
  \item[glibc] the GNU C library of Linux; \code{manylinux\_2\_28} means version 2.28 or newer.
  \item[sdist] source distribution: the sources and \code{pyproject.toml}, compiled on install.
  \item[site-packages] the folder of an environment where packages are installed.
  \item[wheel tag] the \code{python-abi-platform} part of a wheel's name.
\end{description}
